% Part: model-theory
% Chapter: interpolation
% Section: introduction

\documentclass[../../../include/open-logic-section]{subfiles}

\begin{document}

\olfileid{mod}{int}{int}

\olsection{مقدمه}

قضیهٔ درون‌یابی نتیجهٔ زیر است: فرض کنید $\Entails !A \lif !B$.
آنگاه !!a{sentence}~$!C$ وجود دارد چنان‌که
$\Entails !A \lif !C$ و $\Entails !C \lif !B$. افزون بر این، هر
!!{constant}، !!{function} و !!{predicate} (جز~$\eq$) که در~$!C$ رخ
می‌دهد، هم در~$!A$ و هم در~$!B$ رخ می‌دهد. !!{sentence}~$!C$
\emph{درون‌یابِ}~$!A$ و~$!B$ نامیده می‌شود.

قضیهٔ درون‌یابی به‌خودی‌خود جالب است، اما اهمیت اصلی آن در این است که
می‌توان از آن برای اثبات نتایجی دربارهٔ تعریف‌پذیری در یک نظریه و نیز
شرایطی استفاده کرد که تحت آن‌ها ترکیب دو نظریهٔ سازگار نظریه‌ای سازگار
به دست می‌دهد. نتیجهٔ نخست به قضیهٔ تعریف‌پذیری بث و نتیجهٔ دوم به
قضیهٔ سازگاری مشترک رابینسون معروف است.

\end{document}
